\documentclass{standalone}
\usepackage{circuitikz}
\usetikzlibrary{calc}
\definecolor{circuitnotemark}{HTML}{FE00FE}
\begin{document}
\begin{circuitikz}[american, line width=0.8pt]
\ctikzset{bipoles/length=1.2cm}
\foreach \x in {0,2,4,6} {\foreach \y in {0,-2,-4} {\fill[gray, opacity=0.35] (\x,\y) circle (1.1pt);}}
\node[gray, font=\scriptsize] at (0,0.84) {1};
\node[gray, font=\scriptsize] at (2,0.84) {2};
\node[gray, font=\scriptsize] at (4,0.84) {3};
\node[gray, font=\scriptsize] at (6,0.84) {4};
\node[gray, font=\scriptsize, anchor=east] at (-0.84,0) {a};
\node[gray, font=\scriptsize, anchor=east] at (-0.84,-2) {b};
\node[gray, font=\scriptsize, anchor=east] at (-0.84,-4) {c};
\coordinate (b2) at (2,-2);
\coordinate (b1) at (0,-2);
\coordinate (a4) at (6,0);
\coordinate (c4) at (6,-4);
\node[spdt] (part-S1) at (b2) {}; % line 2
\draw (b1) |- (part-S1.in); % line 4
\draw (part-S1.out 1) -| (a4); % line 5
\draw (part-S1.out 2) -| (c4); % line 6
\path (10,-0.226) rectangle (11.861,0.226); % line 8
\node[anchor=west, circuitnotemark, font=\footnotesize] at (10,0) {X}; % line 8
\end{circuitikz}
\end{document}
